\chapter{Comparing the two implementations}
\label{ch:comparing}

Two implementations of one program are an opportunity: whatever differs between them depends
on the platform, and whatever is common does not. This chapter describes three ways to compare
Direct3D and Metal versions of Texture Explorer, from the most informal to the most precise:
reading the code, running both programs, and measuring their results.

\section{Reading the code side by side}
Every source file has a counterpart with the same name, the same functions and the same order
(Table~\ref{tab:pairs}), so the two can be opened next to each other in any editor, or compared
with a difference tool from the root of the repository:
\begin{lstlisting}[language={},numbers=none,frame=none,basicstyle=\ttfamily\small]
git diff --no-index src/Inspector.cpp mac/src/Inspector.cpp
git diff --no-index --word-diff shaders/pbr.hlsl mac/shaders/pbr.metal
\end{lstlisting}
The amount of difference is itself instructive. The geometry and the inspector differ only in
vector types; the shaders only in the syntax of Table~\ref{tab:hlslmsl}; the user interface in a
handful of lines; the renderer and the main program almost everywhere.

\begin{table}[ht]
\centering\small
\begin{tabularx}{\linewidth}{@{}ll>{\raggedright\arraybackslash}X@{}}
\toprule
\textbf{Windows} & \textbf{macOS} & \textbf{What differs}\\
\midrule
--- & \file{mac/src/VecMath.h} & the parts of DirectXMath the program needs (Section~\ref{sec:simd})\\
\file{src/Mesh.*} & \file{mac/src/Mesh.*} & vector types only\\
\file{src/Material.*} & \file{mac/src/Material.*} & Metal textures and views; command buffer for mipmaps\\
\file{src/Inspector.*} & \file{mac/src/Inspector.*} & vector types only\\
\file{shaders/pbr.hlsl} & \file{mac/shaders/pbr.metal} & HLSL syntax $\to$ MSL syntax (Table~\ref{tab:hlslmsl})\\
\file{shaders/gizmo.hlsl} & \file{mac/shaders/gizmo.metal} & idem\\
\file{src/Renderer.*} & \file{mac/src/Renderer.*} & Direct3D~11 $\to$ Metal (Table~\ref{tab:d3dmetal})\\
\file{src/App.*} & \file{mac/src/App.*} & texture handles, paths, opening a URL, Retina scale, Command-R\\
\file{src/main.cpp} & \file{mac/src/main.mm} & Win32 loop $\to$ AppKit delegate and \code{MTKView}\\
\file{CMakeLists.txt} & \file{mac/CMakeLists.txt} & back ends of ImGui, metal-cpp, frameworks, bundle\\
\bottomrule
\end{tabularx}
\caption{The files of the two implementations.}
\label{tab:pairs}
\end{table}

The literate version of the program in \file{LiterateP} makes this comparison precise. It
contains both implementations, and every piece of code that is identical in both is written
once and \emph{used} by both: the woven document says so under the piece, and its index lists
all of them. Since the tool that checks the literate program regenerates both implementations
from it and compares them with the repository byte for byte, the list of shared pieces cannot
become outdated.

\section{Running both programs}
Both programs start with the same material, shape, camera, light and settings, and lay out their
panels the same way; Figures~\ref{fig:pyramid} and~\ref{fig:macpyramid} were taken that way. To
compare them:
\begin{enumerate}
\item Start both programs and choose the same material, shape and view mode.
\item Pin the probe on the same texel. Clicking in the large image of the Inspector is easiest;
the $(u,v)$ shown tells how close the two probes are.
\item Compare the readouts. The values of the maps and the tangent-space normal depend only on
$(u,v)$ and must agree to the digits shown; the world-space vectors depend also on the shape and
on the rotation, which can be set to the same angle with its slider.
\item Compare the images, in the lit view and in the debug views. The ``UV coordinates'',
``World normal'' and ``Tangent-space normal'' views turn the data of the shaders into colours,
so a screenshot of each can be compared pixel by pixel with an image tool.
\end{enumerate}
The images will be almost but not exactly identical, and the reasons are worth discussing:
\begin{itemize}
\item Neither Direct3D nor Metal requires bit-identical results from different GPUs. Texture
filtering, in particular anisotropic filtering, multisample coverage and the precision of
functions such as \code{pow} and \code{normalize} are specified only up to a tolerance.
\item The lines of the arrows are one \emph{pixel} wide, so they look half as thick on a Retina
display.
\item The fonts differ (Segoe~UI and San Francisco), and with them the sizes of the panels in
pixels.
\end{itemize}

\section{Measuring: the comparison program}
\label{sec:probe}
The CPU half of the program---the shapes, the sampling of the maps and the inspector---can be
compared exactly, and \file{mac/compare} contains a program that does it on a Mac. The obstacle
is that the two implementations never run on the same computer. However, the CPU half of the
Windows version needs only standard C++, DirectXMath and a few declarations from Direct3D.
DirectXMath is a header-only library that Microsoft publishes for other platforms too, and the
Direct3D declarations can be replaced by stand-ins that compile and do nothing. So the very files
\file{src/Mesh.cpp}, \file{src/Material.cpp} and \file{src/Inspector.cpp}, unmodified, can be
compiled on a Mac next to their macOS counterparts:
\begin{lstlisting}[language={},numbers=none,frame=none,basicstyle=\ttfamily\small]
cmake -S mac/compare -B mac/compare/build
cmake --build mac/compare/build --target compare
\end{lstlisting}

\paragraph{How it works.} The source file \file{probe.cpp} is compiled twice: once with the
Windows files (and the stand-ins of \file{mac/compare/shims} first in the include path), and once
with the files of the port. Both programs print the same report: the size of each mesh, every
997th vertex, and the full answer of the inspector---colour, bump, displacement, roughness,
slope, $\nts$, and the world-space position, $\Nv$, $\vect n'$, $\Tv$ and $\Bv$---for all ten
materials on all four shapes, at the corners of an $8\times8$ grid of texels (where wrapping
matters) and on a second grid between texel centres (where interpolation matters). The script
\file{compare.py} matches the two reports line by line and prints, for each quantity, the largest
difference. The stand-in for Direct3D is itself instructive: it lists exactly how much of
Direct3D the loading of a material touches.

\excerpt{shimdevice}

\paragraph{Results.} At the time of writing, the 7\,149 lines of the two reports compare as
follows:

\begin{center}\small
\begin{tabular}{@{}lr@{}}
\toprule
\textbf{Quantity} & \textbf{Largest difference}\\
\midrule
mesh sizes and index sums & identical\\
vertex positions, normals and tangents & $2\cdot10^{-7}$\\
vertex texture coordinates & $0$\\
colour, bump, displacement and roughness at the probe & $0$\\
slope of the height field & $0$\\
tangent-space normal $\nts$ & $2\cdot10^{-7}$\\
world-space $\Nv$, $\vect n'$, $\Tv$, $\Bv$ & $2\cdot10^{-7}$\\
world-space position of the probe & $1\cdot10^{-6}$\\
\bottomrule
\end{tabular}
\end{center}

The values read from the maps agree exactly, because both versions sample them with the same
statements (Section~\ref{sec:cpusampling}). The vectors differ in the last one or two bits of a
\code{float}, $2^{-23}\approx1.2\cdot10^{-7}$: DirectXMath and \code{simd} perform the same
operations, but not always in the same way---a normalization may divide by a square root in one
library and multiply by a reciprocal square root in the other. The script accepts differences up
to $10^{-5}$, a hundred times the rounding differences and far below anything visible in the
inspector, and fails otherwise. A real discrepancy is much larger: replacing the factor
$\frac12$ of the central differences by $0.6$ in one version makes the script fail with a
difference of $0.03$ in the slope, and it names the material, shape and texel where the largest
difference occurs.

\section{Exercises}
\begin{enumerate}[label=\textbf{C\arabic*.}]
\item Change the bump normal in one version only, for example to forward differences
(Section~\ref{sec:centraldiff}), and run the comparison. Which quantities does it flag? Why does
the slope change but not the values of the maps?
\item Compare \file{shaders/pbr.hlsl} and \file{mac/shaders/pbr.metal} with a word difference.
Classify every difference into the rows of Table~\ref{tab:hlslmsl}. Is there any difference in the
mathematics?
\item In \file{mac/shaders/pbr.metal}, replace one \code{packed\_float3} of \code{Frame} by
\code{float3}. Predict what happens on the screen, and explain why the \code{static\_assert} of
the C++ structure does not detect the mistake.
\item In the macOS viewport, write \code{view * proj} instead of \code{proj * view}. Describe the
image, and explain why the Windows version writes the product in the other order.
\item On a Retina display, render the scene at the size of the panel in points instead of pixels.
How does the image change, and why?
\item Capture the ``World normal'' view in both programs at the same window size and compute the
difference image. Where are the differences largest, and why there?
\item The comparison program covers only the CPU half. Design a comparison of the GPU half:
what would each program have to write to a file, and what tolerance would be reasonable?
\end{enumerate}
